\documentclass[10pt,a4paper]{amsart}
\usepackage[margin=19mm]{geometry}
\usepackage{amsmath,amssymb,mathtools}
\usepackage[scr=rsfs,cal=euler]{mathalfa}
\usepackage{xfrac}
\usepackage[unicode,hidelinks,bookmarks=false]{hyperref}
\newcommand{\ie}{i.e.\ }
\newcommand{\cf}{cf.\ }
\newcommand{\resp}{respectively\ }
\newcommand{\Subsection}{\subsection}
\providecommand{\nobreakdash}{\nobreak\hbox{-}}
\setlength{\emergencystretch}{2em}
\multlinegap0pt
\usepackage{amssymb}
\usepackage{xcolor}
\usepackage{enumerate}
\newtheorem{lemma}{Lemma}
\title{Sharp Transitions for Localized Solutions to a Diophantine Inequality}
\author{Ataleshvara Bhargava}
\date{Source: arXiv:2608.18350v1 (CC BY 4.0)}
\begin{document}
\maketitle

\noindent\emph{Source and licence.} Sharp Transitions for Localized Solutions to a Diophantine Inequality. Version arXiv:2608.18350v1 (CC BY 4.0), distributed by arXiv under the Creative Commons Attribution 4.0 International licence, http://creativecommons.org/licenses/by/4.0/, as recorded in the arXiv licence metadata for this exact version and shown on the abstract page https://arxiv.org/abs/2608.18350v1. Full licence notice: \texttt{licenses/arxiv-2608.18350-CC-BY-4.0.txt}.\par\smallskip

\noindent\textit{Editorial scope.} Counted unit: the article's lemma lemma (source0.tex lines 141--142). The article's complete original proof of it is delivered with it (source0.tex lines 144--203, one \texttt{proof} environment of 60 lines). No mathematics is written, repaired or re-derived: the statement and the proof are the article's own lines, in the article's own notation, and no retained line is substituted or re-flowed.  Retained uncounted as the source-local context the proof needs: lines 91--93.  Nothing was omitted from any retained span, so the package carries no omission record.  No retained line contains a citation, so the wrapper carries no bibliography and no bibliography entry.  No retained line contains a figure environment, an image-inclusion command or drawing code, so no image is delivered and the package declares no asset dependency.  Editorial formatting only, in addition: an AMS \texttt{amsart} preamble that restates the article's own declarations and macros used by the retained lines, namely the environments \texttt{lemma} on separate counters, together with the standard packages those lines need.  The retained environments print in this document as Lemma 1 in order of appearance, whereas the article prints the same items with the numbering of its own full document.  The complete native source of the article as distributed in the arXiv e-print for this exact version is bundled unchanged as \texttt{sources/207985/source0.tex}.
\begin{align} \label{eq:main}
|x_1^{\theta} + \cdots +x_s^{\theta} - R| < \tau, \quad X-Y < x_i \leq X+Y \quad \text{ for each } \ i = 1, \ldots, s.
\end{align}

\begin{lemma} Fix $\tau > 0$ and $\theta > 2$ with $\theta \notin \mathbb{N}$, and let $s \in \mathbb{N}$ with $s \geq 2$. Let $M \in \mathbb{N}$ be sufficiently large. Fix $c < (\lfloor s/2 \rfloor (\theta-1))^{-1/2} $. Then there are no solutions to \eqref{eq:main} when $R = sM^{\theta}+s\theta M^{\theta-1}$. 
\end{lemma} 

\begin{proof} Let $R = sM^{\theta}+s\theta M^{\theta-1}$, where $M$ is a large integer. Recall that $X = (R/s)^{1/\theta}$. Assume for sake of contradiction that there are $x_1,\ldots, x_s$ such that \eqref{eq:main} is true, with $Y = c\sqrt{X}$. Let $A_i = x_i-M$ for each $i$. Note that since $X^{\theta}= R/s = M^{\theta}+\theta M^{\theta-1}$, we have $X \geq M$. We may therefore deduce from the Mean Value Theorem that 
\[ \theta M^{\theta-1}|X-M| \leq |X^{\theta}-M^{\theta}| = \theta M^{\theta-1}. \]
Therefore we have $M \leq X \leq M+1$, and also 
\begin{align*} 
|A_i| & = |x_i-M| \leq |x_i-X|+|X-M| \leq  cX^{1/2} +1 \leq c(M+1)^{1/2}+1.
\end{align*}
In particular, for any fixed $\delta > 0$ to be chosen later we have $|A_i| \leq c(1+\delta)M^{1/2}$ for all large $M$. 
One more observation about $A_i$ is that we cannot have $x_i = M$ for all $i$ and for large $M$. This would mean $x_1^{\theta}+\cdots +x_s^{\theta} = sM^{\theta}$. However, $sM^{\theta}$ is not within $\tau$ of $R$, and so this contradicts the assumption that $(x_1,\ldots,x_s)$ is a solution to \eqref{eq:main}. Therefore, we must have 
\begin{align} \label{eq:A_i}
\sum_{i=1}^s A_i^2 = \sum_{i=1}^s |x_i-M|^2 \geq 1. 
\end{align}
Using the Taylor series for the function $h(x) = (M+x)^{\theta}$ centered at the origin with number of terms $k$ and remainder $r_k$ we then have 
\begin{equation} \label{eq:comp}
\begin{aligned}
\tau & > \Big|\sum_{i=1}^s (M+A_i)^{\theta}-sM^{\theta}-s\theta M^{\theta-1}\Big| = \Big|\sum_{i=1}^s \sum_{j=0}^k \binom{\theta}{j}A_i^{j}M^{\theta-j}+\sum_{i=1}^s r_k(A_i)-sM^{\theta}-s\theta M^{\theta-1}\Big| 
\\ & = \Big|\theta M^{\theta-1}\big(-s+\sum_{i=1}^s A_i\big)+\sum_{i=1}^s \sum_{j=2}^k \binom{\theta}{j}A_i^jM^{\theta-j}+\sum_{i=1}^s r_k(A_i)\Big| 
\\ & \geq \Big|\theta M^{\theta-1}\big(-s+\sum_{i=1}^s A_i\big)\Big|-\Big|\sum_{i=1}^s \sum_{j=2}^k \binom{\theta}{j}A_i^jM^{\theta-j}+\sum_{i=1}^s r_k(A_i)\Big|. 
\end{aligned}
\end{equation}
Now, the remainder is given as 
\[r_k(A_i) = \binom{\theta}{k+1}(M+z)^{\theta-k-1}A_i^{k+1}. \] Since $|A_i| \leq c(1+\delta)M^{1/2}$, one has $M+A_i = M +O(M^{1/2})$, and so 
\[ |r_k(A_i)| \ll_{\theta,k,c,\delta} M^{\theta-k-1}M^{(k+1)/2} \ll M^{\theta-(k+1)/2}. \] 
When $k \geq 2$ it follows that 
\[ \Big|-s+\sum_{i=1}^s A_i\Big| \leq \theta^{-1}\Big|\sum_{i=1}^s \sum_{j=2}^k \binom{\theta}{j}A_i^jM^{1-j}\Big|+O_{\tau, \theta, k, c, \delta}(M^{-1/2}) \]
\[ \leq \theta^{-1}\sum_{i=1}^s \sum_{j=2}^k \Big|\binom{\theta}{j}\Big|(c(1+\delta)M^{1/2})^jM^{1-j}+o(1) \leq s\theta^{-1} \sum_{j=2}^k \Big| \binom{\theta}{j} \Big| (c(1+\delta))^j M^{1-j/2}+o(1). \]
Noting that $M^{1-j/2} = O(M^{-1/2})$ for all $j > 2$, for any $\delta > 0$ we have 
\begin{align} \label{eq:sA_i}
\Big|s-\sum_{i=1}^s A_i\Big| \leq \frac{1}{2}sc^2(1+\delta)^2(\theta-1)+o_{\tau, \theta, k, c, \delta}(1). 
\end{align}
We now consider two cases, according to whether $s$ is even or odd. Suppose first that $s$ is even. Let $k = 4$. 
Since $c < \sqrt{2}(s(\theta-1))^{-1/2}$ we may find $\delta > 0$ such that $2^{-1}s(c(1+\delta))^2(\theta-1) < 1$. Therefore, for all large $M$ we see that the inequality 
\[ \Big|-s+\sum_{i=1}^s A_i\Big| < 1 \]
holds. Since the left hand side is an integer, we deduce that this integer must be $0$. Substituting this into the second line of \eqref{eq:comp} and using $k = 4$, one has 
\[ \tau > \Big|\sum_{i=1}^s \sum_{j=2}^4 \binom{\theta}{j} A_i^jM^{\theta-j} +\sum_{i=1}^s r_4(A_i)\Big| \geq \Big|\sum_{i=1}^s \binom{\theta}{2}M^{\theta-2}A_i^2\Big| - \Big|\sum_{i=1}^s \sum_{j=3}^4 \binom{\theta}{j}M^{\theta-j}A_i^j+\sum_{i=1}^s r_4(A_i)\Big|. \]
Rearranging and using the bounds $r_4(A_i) \ll M^{\theta-5/2}$ and $A_i \ll M^{1/2}$ we may see that 
\[ \binom{\theta}{2}M^{\theta-2}\sum_{i=1}^s A_i^2 \ll \tau + \sum_{j=3}^4 M^{\theta-j/2-1}\sum_{i=1}^s A_i^2+M^{\theta-5/2} \ll 1 +M^{\theta-5/2}\sum_{i=1}^s A_i^2+M^{\theta-5/2} . \]
Since \eqref{eq:A_i} holds and $\theta > 2$, this is impossible, since the terms on the right are too small to majorize the term on the left.  

Now instead suppose $s$ is odd. Since $c < (\lfloor s/2 \rfloor(\theta-1))^{-1/2}$ we may choose $\delta > 0$ such that $c < (\lfloor s/2 \rfloor(\theta-1)(1+\delta)^3)^{-1/2}$. Since $s$ is odd, $\lfloor s/2 \rfloor = (s-1)/2$ and by \eqref{eq:sA_i} we must have 
\[ \Big|s-\sum_{i=1}^s A_i\Big| \leq \frac{1}{2}sc^2(1+\delta)^2(\theta-1)+o(1) < \frac{s}{(1+\delta)(s-1)}+o(1) < 2+o(1). \]
Define $\chi$ to be the expression within the absolute values on the left hand side. For sufficiently large $M$ it must then be that $\chi \in \{ -1, 0 , 1\}$. If $\chi = 0$, then we may argue exactly as in the previous case to derive the same contradiction. If instead $\chi = \pm 1$, then inserting this into \eqref{eq:comp} with $k = 2$ and using the Triangle inequality shows 
\[ |\theta M^{\theta-1}\chi| < \tau +\Big| \binom{\theta}{2} M^{\theta-2} \sum_{i=1}^s A_i^2\Big|+O(M^{\theta-3/2}), \]
whence 
\begin{align} \label{eq:1bound}
1 < \frac{\theta-1}{2M}\sum_{i=1}^sA_i^2+O(M^{-1/2})
\end{align}
Since $s$ is odd, one of the two sets $\{1 \leq i \leq s: A_i > 0\}$ and $\{1 \leq i \leq s :A_i \leq 0\}$ contains at most $(s-1)/2$ elements. Let $U$ be the smaller set, and let $V$ be the larger. It follows that 
\[ \sum_{i \in U} |A_i| \leq 2^{-1}c(s-1)(1+\delta)M^{1/2} \]
since $|A_i| \leq c(1+\delta)M^{1/2}$ for all $i$. Note that all $A_i$ with $i \in V$ have the same sign. Hence, by definition of $\chi$ one has 
\[ \sum_{i \in V} |A_i| = \big|\sum_{i \in V} A_i\big| = \big|s-\chi-\sum_{i \in U} A_i\big| \leq s+1+2^{-1}c(s-1)(1+\delta)M^{1/2}. \] 
Moreover, we have 
\[ \sum_{i=1}^s |A_i| = \sum_{i \in U} |A_i| +\sum_{i \in V} |A_i| \leq c(s-1)(1+\delta)M^{1/2}+s+1 \leq c(s-1)(1+\delta)^2M^{1/2} \]
for all sufficiently large $M$. Therefore, 
\[ \sum_{i=1}^s |A_i|^2 \leq c(1+\delta)M^{1/2}\sum_{i=1}^s |A_i| \leq c^2(s-1)(1+\delta)^3M. \]
Using this in \eqref{eq:1bound} we deduce 
\[ 1 < \frac{1}{2}c^2(\theta-1)(s-1)(1+\delta)^3+O(M^{-1/2}). \]
However, the upper bound on $c$ implies that 
\[ \frac{1}{2}c^2(\theta-1)(s-1)(1+\delta)^3 < 1, \]
and so the above inequality is impossible if $M$ is too large, which is a contradiction. 
\end{proof} 

\end{document}
